% TOP_PROBLEM: {"problemId": "problem.top500-omission-w045049-positive-type-a-schubert-counting-rule", "canonicalTitle": "Positive Counting Rule for Type-A Schubert Structure Constants", "releaseRank": 490, "primaryDomain": "algebra_representation_category", "primaryDomainLabel": "Algebra, representation & category theory", "displayStatus": "open_with_solved_subcases", "releaseStatus": "open_with_solved_subcases"}
% !TeX program = lualatex
% Definition and sources only; this file is not an LLM solution attempt.
\documentclass[11pt]{article}
\usepackage[margin=25mm]{geometry}
\usepackage{fontspec}
\setmainfont{Latin Modern Roman}
\usepackage{amsmath,amssymb,mathtools}
\usepackage{unicode-math}
\setmathfont{Latin Modern Math}
\usepackage{xurl,hyperref}
\hypersetup{colorlinks=true,urlcolor=blue}
\setcounter{secnumdepth}{0}
\setlength{\parindent}{0pt}
\setlength{\parskip}{5pt}
\setlength{\emergencystretch}{3em}
\begin{document}
\section*{\texorpdfstring{Positive Counting Rule for Type-A Schubert Structure Constants}{Positive Counting Rule for Type-A Schubert Structure Constants}}\label{490}
\subsection{Definitions and mathematical statement}
Let $\mathfrak S_w$ be the ordinary type-A Schubert polynomial indexed by a finite permutation $w$, in the stable polynomial convention. Its nonnegative integer multiplication constants are defined by
\[
 \mathfrak S_u\mathfrak S_v=\sum_w c_{u,v}^{w}\mathfrak S_w.
\]
Encode $u,v,w$ explicitly in one-line notation as a bit string $x$. The desired positive counting rule has the precise complexity meaning
\[
 \begin{gathered}
 \exists\text{ polynomial }p\ \exists\text{ polynomial-time predicate }R\\
 \forall x=(u,v,w):\quad
 c_{u,v}^{w}=\#\{y\in\{0,1\}^{p(|x|)}:R(x,y)=1\}.
 \end{gathered}
\]
Thus the function belongs to $\#\mathrm P$ under this encoding. A signed sum with cancellations or a nonconstructive positivity proof is not such a rule. The predicate and polynomial must be uniform across all permutations; there is no fixed-rank restriction.
\par\smallskip\textit{Formulation and concept references:} \href{https://arxiv.org/html/2412.18984}{[S1]}.

\subsection{Short English statement}
Can every ordinary type-A Schubert multiplication coefficient be counted by a single polynomial-time-verifiable family of positive witnesses?
\subsection{Sources}
\begin{itemize}
\item[S1]Igor Pak and Colleen Robichaux, Positivity of Schubert Coefficients (2024 preprint; Comptes Rendus Mathématique 363 (2025), 1507–1515).
\url{https://arxiv.org/html/2412.18984}.
\textit{Formulation source.}
\end{itemize}
\end{document}
